\section{The whole module on one page}\label{sec:map}

\subsection{The big picture: one story, five chapters}
The whole module is about one object, the \textbf{$(A,B,C,D)$ linear system}
\[
\boxed{\;\frac{dX}{dt}=AX+BU,\qquad Y=CX+DU\;}
\]
with input $U(t)$, state $X(t)$ and output $Y(t)$. Every chapter looks at it from a different angle:

\begin{center}
\renewcommand{\arraystretch}{1.3}
\begin{tabular}{@{}p{3.3cm}p{5.6cm}p{6.2cm}@{}}
\toprule
\textbf{Chapter} & \textbf{Question it answers} & \textbf{Main tools}\\\midrule
1. Description \fr{6--30} & How do I write a machine or an ODE as a system? & block diagrams, reduction of order, companion matrix\\
2. Matrix theory \fr{31--85} & How do I \emph{solve} the system? & $\det$, eigenvalues, $\exp(tA)$, solution formula, $T(s)=D+C(sI-A)^{-1}B$, realisation\\
3. Laplace \fr{86--134} & How do I turn calculus into algebra? & table, properties, convolution, partial fractions, Jordan form\\
4. Transfer functions \& stability \fr{135--161} & Does the output blow up? How does it respond to frequency $\omega$? & gain, phase, Nyquist plot, exponential/marginal/BIBO stability\\
5. Feedback control \fr{162--218} & How do I \emph{make} a system stable? & stable polynomials, $\mathcal S$, Nyquist criterion, PID, coprime factorisation, Youla, Lyapunov, Sylvester\\
\bottomrule
\end{tabular}
\end{center}

\subsection{The ``chain'' you must be able to travel along}
\begin{center}
\begin{tikzpicture}[node distance=6mm and 9mm, every node/.style={draw,rounded corners,align=center,font=\small,fill=teal!5}]
\node (ode) {ODE / integral equation\\ $y''+\beta y'+\gamma y=u$};
\node (bd) [right=of ode] {block diagram\\(integrators, no $d/dt$)};
\node (abcd) [right=of bd] {$(A,B,C,D)$\\$\dot X=AX+BU$};
\node (T) [below=of abcd] {transfer function\\$T(s)=D+C(sI-A)^{-1}B$};
\node (st) [left=of T] {poles $=$ eigenvalues\\ stability (BIBO)};
\node (fb) [left=of st] {feedback $\dfrac{G}{1+GK}$\\ stabilise};
\draw[-Stealth] (ode)--(bd); \draw[-Stealth] (bd)--(abcd); \draw[-Stealth] (abcd)--node[draw=none,fill=none,right]{$\Lap$}(T);
\draw[-Stealth] (T)--(st); \draw[-Stealth] (st)--(fb);
\draw[-Stealth,dashed] (T.east) to[bend right=40] node[draw=none,fill=none,right]{realise} (abcd.east);
\end{tikzpicture}
\end{center}

\subsection{Question-type finder (use this in the exam)}
\begin{center}
\renewcommand{\arraystretch}{1.25}
\begin{tabular}{@{}p{7.3cm}p{8cm}@{}}
\toprule
\textbf{If the question says\dots} & \textbf{Go to / do this}\\\midrule
``Express as a block diagram (no differentiators)'' & Integrate both sides until no derivatives are left; \S\ref{sec:bd}\\
``Write in matrix form $(A,B,C,D)$'' / ``introduce state $v=dx/dt$'' & Reduction of order, \S\ref{sec:reduction}\\
``Find a matrix with characteristic polynomial $p$'' & Companion matrix, \S\ref{sec:companion}\\
``Find $(sI-A)^{-1}$'' & Adjugate / triangular trick, \S\ref{sec:resolvent}\\
``Find $\exp(tA)$'' & Diagonalise, Jordan or inverse Laplace of $(sI-A)^{-1}$, \S\ref{sec:exp}\\
``Solve $\dot X=AX+BU$, $X(0)=X_0$'' & Variation of constants, \S\ref{sec:solution}\\
``Compute the transfer function'' & $T(s)=D+C(sI-A)^{-1}B$, \S\ref{sec:tf}\\
``Realise $T(s)$ as a SISO'' & Companion realisation, \S\ref{sec:realise}\\
``Rank of $[C;CA;\dots;CA^{n-1}]$'' & \S\ref{sec:rank}\\
``Laplace transform of \dots{} from the definition'' & Integrate by parts / exponential trick, \S\ref{sec:lapdef}\\
``Solve the IVP/integral equation by Laplace'' & \S\ref{sec:ivp}, \S\ref{sec:conv}\\
``Jordan form, minimal polynomial, multiplicities'' & \S\ref{sec:jordan}\\
``Gain and phase / Bode / Nyquist'' & \S\ref{sec:freq}\\
``Is the system stable / BIBO stable?'' & Eigenvalues of $A$ / poles of $T$, \S\ref{sec:stab}\\
``Is the polynomial stable?'' / ``for which $k$ is it stable?'' & \S\ref{sec:stabpoly}\\
``Find a controller $-a-b/s$ / PID that stabilises $G$'' & \S\ref{sec:pid}\\
``Coprime factorisation / stabilising controller'' & \S\ref{sec:coprime}\\
``Internally stable / well-posed'' & \S\ref{sec:internal}\\
``Positive definite / Lyapunov / Sylvester'' & \S\ref{sec:posdef}--\S\ref{sec:sylv}\\
``Prove that \dots'' & the grey \emph{Proof} boxes; list in \S\ref{sec:proofs}\\
\bottomrule
\end{tabular}
\end{center}

\begin{exambox}[title={Exam technique for this module}]
\begin{enumerate}
\item \textbf{State the formula first}, then substitute. E.g.\ write ``$T(s)=D+C(sI-A)^{-1}B$'' before doing any algebra. This earns method marks even if the arithmetic slips.
\item \textbf{Keep track of shapes.} $A$ is $n\times n$, $B$ is $n\times k$, $C$ is $m\times n$, $D$ is $m\times k$. If shapes do not match, you have made an error.
\item \textbf{Check your answer cheaply:} $\det$ = product of eigenvalues, $\tr$ = sum of eigenvalues; put $t=0$ into $\exp(tA)$ (should give $I$); put $s\to\infty$ into $T(s)$ (should give $D$); put $t=0$ into an IVP solution.
\item \textbf{Use the lecture notes' sign conventions}: $\chi_A(s)=\det(sI-A)$ (not $\det(A-sI)$); the adjoint $A'$ is the conjugate transpose; LHP is the \emph{open} left half plane.
\item \textbf{Say why}: ``all poles have negative real part, hence $T\in\mathcal S$'' is a full-mark sentence; ``it is stable'' is not.
\end{enumerate}
\end{exambox}
\newpage
